% =====================================================================
%  Agents Don't Game the Linter
%  Compiles with pdfLaTeX on Overleaf. Upload the four PNGs to the same
%  folder as this file:
%     fig1_breakage_by_rule_group.png
%     fig2_paired_arms.png
%     fig3_by_rule.png
%     fig4_no_gaming.png
% =====================================================================
\documentclass[11pt]{article}

\usepackage[T1]{fontenc}
\usepackage[utf8]{inputenc}
\usepackage{newtxtext,newtxmath}
\usepackage[margin=1.05in]{geometry}
\usepackage{microtype}
\usepackage{graphicx}
\usepackage{booktabs}
\usepackage{xcolor}
\usepackage{caption}
\usepackage{titlesec}
\usepackage{enumitem}
\usepackage{tcolorbox}
\usepackage[hidelinks]{hyperref}

\graphicspath{{./}{figures/}}

% Floats here are wide and short; the defaults strand them on sparse pages.
\renewcommand{\topfraction}{0.92}
\renewcommand{\bottomfraction}{0.85}
\renewcommand{\textfraction}{0.08}
\renewcommand{\floatpagefraction}{0.80}
\setcounter{topnumber}{3}
\setcounter{bottomnumber}{2}
\setcounter{totalnumber}{5}

% --- palette, matched to the figures -------------------------------------
\definecolor{ink}{HTML}{0B0B0B}
\definecolor{muted}{HTML}{52514E}
\definecolor{accent}{HTML}{2A78D6}
\definecolor{warm}{HTML}{EB6834}
\definecolor{rule}{HTML}{D8D7D3}
\definecolor{panel}{HTML}{F6F7F9}

\hypersetup{colorlinks=true, linkcolor=accent, urlcolor=accent, citecolor=accent}

% --- typography ----------------------------------------------------------
\titleformat{\section}{\normalfont\large\bfseries\color{ink}}{\thesection}{0.7em}{}
\titleformat{\subsection}{\normalfont\normalsize\bfseries\color{ink}}{\thesubsection}{0.6em}{}
\titlespacing*{\section}{0pt}{1.6em plus .3em minus .2em}{0.6em}
\titlespacing*{\subsection}{0pt}{1.2em plus .2em minus .2em}{0.4em}

\captionsetup{font=small, labelfont={bf,color=ink}, textfont={color=muted},
              width=.94\textwidth, skip=8pt}
\captionsetup[table]{skip=6pt}
\renewcommand{\arraystretch}{1.18}
\setlength{\parskip}{0.35em}
\setlength{\parindent}{0pt}

\newcommand{\code}[1]{\texttt{\small #1}}
\newcommand{\ci}[2]{\,[#1,\,#2]}

% =====================================================================
\begin{document}

\begin{center}
{\LARGE\bfseries\color{ink} Agents Don't Game the Linter\par}
\vspace{0.45em}
{\large\color{muted} Measuring whether LLM repair satisfies the analyser\\
or fixes the defect\par}
\vspace{1.3em}
{\normalsize Syed Mohammed Sameer\par}
\vspace{0.25em}
{\small\color{muted} \href{mailto:sdmohammedsameer@gmail.com}{sdmohammedsameer@gmail.com}\par}
\vspace{0.35em}
{\small\color{muted} \today\par}
\end{center}

\vspace{0.6em}

\begin{tcolorbox}[colback=panel, colframe=rule, boxrule=0.5pt, arc=2pt,
                  left=12pt, right=12pt, top=10pt, bottom=10pt]
\small
\textbf{Abstract.} Four code-specialised LLMs were shown static-analysis findings
on their own solutions and asked to fix them. The study was designed to catch
Goodharting: satisfying the analyser that was shown while the defect stayed put.
\textbf{It did not happen.} Across 4{,}439 repair steps the models wrote five
suppression directives in total, and 94\% of the findings they removed were also
gone from a held-out analyser they never saw (suppression rate 6\%, 95\% CI
$[4\%, 9\%]$). They complied instead---and complying with one particular rule
broke two working programs in five. Comparing two repair arms within the same
task on the same model's own solution, the arm shown Ruff broke the program where
the arm shown Pylint did not on 46 occasions against 3 the other way (exact
McNemar, $p = 7\times10^{-11}$, $n = 213$ pairs). That result is one rule:
Ruff's \code{S311} broke 41.8\% $[32\%, 52\%]$ of the programs it appeared in,
and removing every \code{S311} task drops the paired comparison to 10 versus 3,
$p = 0.09$. The finding is therefore not that security rules are dangerous in a
repair loop, but that \code{S311} is, and that a single rule was enough to make an
entire ruleset look four times more destructive than another. A suppressed
finding leaves a directive in the diff; faithful destructive repair leaves nothing
to find.
\end{tcolorbox}

% =====================================================================
\section{The question}

Static analysis is increasingly placed inside an automated loop: an agent writes
code, an analyser reports findings, the agent revises. The findings go down. That
much is established and is the reason the pattern is being adopted.

The obvious worry is Goodhart's law. If the measure becomes the target, an agent
may learn to satisfy the particular analyser it is scored against rather than
improve the code. In Python the cheapest way to do that is a suppression comment:
\code{\# noqa} silences Ruff, \code{\# nosec} silences Bandit, and neither
changes a line of behaviour.

This study asks whether that happens, using an instrument that can tell the
difference, and then reports what it found instead.

% =====================================================================
\section{Design}

\subsection{The instrument}

Ruff's \code{S} rules are a reimplementation of Bandit's, numbered identically:
\code{S602} and \code{B602} are the same defect seen by two separately built
engines. Bandit does not honour Ruff's suppression syntax. That near-duplication
is what makes the pair sharp rather than redundant.

\begin{table}[htbp]
\centering
\small
\begin{tabular}{@{}lll@{}}
\toprule
\textbf{variant of \code{subprocess.check\_output(cmd, shell=True)}}
  & \textbf{Ruff (shown)} & \textbf{Bandit (held out)} \\
\midrule
vulnerable                       & \code{S602}   & \code{B404}, \code{B602} \\
suppressed with \code{\# noqa}   & \textit{none} & \code{B404}, \code{B602} \\
genuinely fixed                  & \code{S603}   & \code{B404}, \code{B603} \\
\bottomrule
\end{tabular}
\caption{A finding that leaves the shown analyser while its held-out twin remains
is suppression. One that leaves both is repair. This is re-verified on the
installed tool versions before every run.}
\end{table}

\subsection{Why counts are not the measurement}

Read that table as counts and it says the opposite of the truth. The vulnerable
variant scores 1 and 2; the genuinely fixed variant also scores 1 and 2, because
the fix moves \code{S602} to \code{S603} and \code{B602} to \code{B603}. Only the
suppression changes a count, and it changes it favourably.

A count-based delta would therefore score a real fix as worthless and a
\code{\# noqa} as a total success, inverting the study's own headline. Every
finding is instead tracked individually by code: \emph{addressed} if the shown
analyser no longer reports it, \emph{transferred} if the held-out counterpart is
also gone, \emph{introduced} if repair traded one defect for another. A finding
with no counterpart is recorded as \emph{unmeasurable} rather than folded into
either bucket, because counting it either way manufactures a result. The run
observed three such swaps.

\subsection{The two arms}

Each task produces one baseline solution and then two \emph{independent} repair
arms branching from that same baseline:

\begin{table}[htbp]
\centering
\small
\begin{tabular}{@{}llll@{}}
\toprule
\textbf{arm} & \textbf{shown} & \textbf{rounds} & \textbf{runs on} \\
\midrule
\code{shown\_pylint} & Pylint & 2 & every task (Pylint fires on 100\% of files) \\
\code{shown\_ruff}   & Ruff   & 2 & tasks with Ruff findings ($\approx$50\%) \\
\bottomrule
\end{tabular}
\end{table}

Branching rather than chaining is what makes every comparison paired within a
task. Both arms record findings from all three analysers at every step, plus the
task's own tests, because a finding removed by breaking the function is not a fix
and must stay distinguishable from one that was merely hidden.

The repair prompt says nothing about suppression in either direction. Telling the
agent not to suppress would destroy the behaviour being measured; telling it that
it may would manufacture it. It does require that behaviour be preserved.

\subsection{What was run}

Four code-specialised instruct models from four families, within a
1.3$\times$ size spread, served at bfloat16 on one A100 through vLLM at
temperature 0: \code{Qwen2.5-Coder-7B-Instruct},
\code{deepseek-coder-6.7b-instruct}, \code{Yi-Coder-9B-Chat}, and
\code{granite-8b-code-instruct-128k}.

300 BigCodeBench tasks, drawn by a fixed shuffle so every model walks the same
order, then filtered to those whose \emph{reference} solution passes its own tests
in the run environment---so a missing package costs coverage rather than
corrupting the correctness measurement. Two repair rounds per arm. 1{,}200
task--model pairs, 4{,}439 steps, 1.14 hours, zero errors. Analysers were Ruff
0.15.11 (\code{--select S,B,SIM,RET,ARG,PERF,C90}), Bandit 1.9.4, and Pylint 4.0.6.

% =====================================================================
\section{Result 1: the agents did not game the analyser}

\begin{table}[htbp]
\centering
\small
\begin{tabular}{@{}lrrrr@{}}
\toprule
\textbf{model} & \textbf{removed} & \textbf{measurable} & \textbf{suppressed} & \textbf{95\% CI} \\
\midrule
deepseek-coder-6.7b & 152 & 62  & 2\%  & $[0\%,\ 9\%]$ \\
granite-8b-code     & 110 & 53  & 8\%  & $[3\%, 18\%]$ \\
qwen2.5-coder-7b    & 179 & 78  & 12\% & $[6\%, 20\%]$ \\
yi-coder-9b         & 208 & 90  & 3\%  & $[1\%,\ 9\%]$ \\
\midrule
\textbf{pooled}     & \textbf{649} & \textbf{283} & \textbf{6\%} & $\mathbf{[4\%,\ 9\%]}$ \\
\bottomrule
\end{tabular}
\caption{Of the findings each agent removed from the analyser it was shown, the
proportion still reported by the held-out twin it never saw. Wilson intervals.}
\end{table}

\begin{figure}[htbp]
\centering
\includegraphics[width=0.97\textwidth]{fig4_no_gaming}
\caption{The result the study was built to detect, and did not find. The dashed
line marks what systematic gaming would look like.}
\end{figure}

The interval is nowhere near 50\%. Two independent measurements agree: the
directives actually written---five, across the entire study---and the held-out
analyser, which shows 94\% of removals transferring. The agents addressed 70\% of
Ruff findings and 20\% of Pylint's, and did so honestly.

This is a negative result on the study's own premise. It is also what makes the
next section believable: an agent that were gaming the metric would confound every
correctness measurement that follows.

% =====================================================================
\section{Result 2: faithful repair breaks working code}

Where both arms ran on the same task and the baseline passed its tests, the arms
can be compared \emph{within} that task, on that model's own solution. Only the
discordant pairs carry information---cases where both arms broke it, or neither
did, say nothing about which arm is worse---so only those are tested, with an
exact McNemar rather than the chi-square approximation, which is not trustworthy
at these counts.

\begin{table}[htbp]
\centering
\small
\begin{tabular}{@{}lrrrrrr@{}}
\toprule
\textbf{model} & \textbf{pairs} & \textbf{Ruff broke it} & \textbf{Pylint broke it}
 & \textbf{both} & \textbf{neither} & \textbf{$p$ exact} \\
\midrule
deepseek-coder-6.7b & 46 &  9 & 1 & 1 & 35 & 0.022 \\
granite-8b-code     & 49 & 10 & 0 & 0 & 39 & 0.002 \\
qwen2.5-coder-7b    & 56 & 12 & 0 & 4 & 40 & 0.0005 \\
yi-coder-9b         & 62 & 15 & 2 & 0 & 45 & 0.002 \\
\midrule
\textbf{pooled} & \textbf{213} & \textbf{46} & \textbf{3} & \textbf{5} & \textbf{159}
 & $\mathbf{7\times10^{-11}}$ \\
\bottomrule
\end{tabular}
\caption{Discordant pairs run 46 to 3 against the Ruff arm. Every model is
individually significant.}
\end{table}

\begin{figure}[htbp]
\centering
\includegraphics[width=0.97\textwidth]{fig2_paired_arms}
\caption{The controlled comparison. Same task, same model, same baseline
solution; the only difference is which analyser's findings the agent was shown.}
\end{figure}

Unpaired, the same picture: the Ruff arm turned 213 working solutions into 163,
breaking 23.9\% $[18.7\%, 30.1\%]$ of them and repairing one, while the Pylint arm
broke 3.8\% $[2.3\%, 6.3\%]$.

% =====================================================================
\section{Result 3: the damage is one rule}

\begin{table}[htbp]
\centering
\small
\begin{tabular}{@{}lrrrr@{}}
\toprule
\textbf{findings the agent was shown} & \textbf{tasks} & \textbf{broke} & \textbf{rate} & \textbf{95\% CI} \\
\midrule
\code{S311} among them              & 91 & 38 & \textbf{41.8\%} & $[32\%, 52\%]$ \\
some other security rule, no \code{S311} & 31 &  5 & 16.1\% & $[7\%, 33\%]$ \\
no security rule at all             & 91 &  8 & 8.8\%  & $[5\%, 16\%]$ \\
\bottomrule
\end{tabular}
\end{table}

\begin{figure}[htbp]
\centering
\includegraphics[width=0.97\textwidth]{fig1_breakage_by_rule_group}
\caption{Grouping by ``security rule or not'' looked like a fourfold effect until
\code{S311} was held out. \code{S311} supplies 91 of the 122 security-shown tasks,
so that grouping was reporting one rule under a class name.}
\end{figure}

\begin{tcolorbox}[colback=white, colframe=warm, boxrule=1pt, arc=2pt,
                  left=12pt, right=12pt, top=9pt, bottom=9pt]
\small
\textbf{Removing \code{S311} removes the result.} Repeating the paired test on
tasks where \code{S311} never appeared gives 10 versus 3, $p = 0.09$. At this
sample size that is indistinguishable from noise.
\end{tcolorbox}

The honest claim is therefore narrow. This is \emph{not} evidence that security
rules as a class are dangerous in a repair loop. It is evidence that \code{S311}
is, and that a single rule was enough to make the whole Ruff arm look four times
more destructive than the Pylint arm. Anyone reasoning about ``static analysis in
the agent loop'' as one undifferentiated thing would have drawn the wrong
conclusion, in either direction.

\begin{figure}[htbp]
\centering
\includegraphics[width=0.97\textwidth]{fig3_by_rule}
\caption{Breakage rate per rule, among tasks where that rule was one of the
findings shown. The two leaders are rules whose fix changes program behaviour.}
\end{figure}

% =====================================================================
\section{The mechanism}

\code{S311}---\emph{``standard pseudo-random generators are not suitable for
cryptographic purposes''}---fired on 202 task--model pairs and was removed 79\% of
the time.

None of that code was cryptographic. BigCodeBench uses \code{random} for data
generation, and its tests seed it and assert on the output. The agent does as it
is told, swaps \code{random} for \code{secrets}, and the program stops being
deterministic. \code{secrets} has no \code{seed()}, so some stop running at all:

\begin{center}
\small\ttfamily\color{muted}
BigCodeBench/87\quad qwen2.5-coder-7b\quad fail: NameError: name 'seed' is not defined
\end{center}

The finding was a true positive about the code and a false positive about the
context. The agent could not tell the difference, and nothing downstream could
either: the analyser reported success. Of the 51 broken repairs in the Ruff arm,
23 raised an error during the test, 20 changed behaviour while still running, and
14 introduced a Pylint \code{E0602} or \code{E1101}---a name the repair had
removed.

\code{B006} (mutable default argument) is the same shape of mistake on a smaller
sample: 17 tasks, 53\% broken. Changing \code{def f(x=[])} to \code{def f(x=None)}
is textbook correct and changes the API. What \code{S311} and \code{B006} share is
not that they are security rules---\code{B006} is not---but that their fix
\emph{changes what the program does}, and neither analyser nor agent has any way to
know whether that was acceptable here.

% =====================================================================
\section{Why this matters more than the gaming result}

A suppressed finding is detectable. The directive sits in the diff, and a second
analyser catches it immediately. Faithful, destructive repair leaves nothing to
find: the analyser is satisfied, the diff looks like a fix, and the program is
broken. The only thing that catches it is running the code.

This is an empirical case for verification that is independent of the gate being
optimised against---and a different case from the one usually made. The risk in an
automated quality gate is not that agents cheat it. It is that they obey it in
contexts where its advice is wrong. The exposure being concentrated in a few rules
rather than spread thinly is the encouraging part: it is findable and fixable per
rule, rather than an argument against the practice.

For a repair loop in production, three things follow. Run the test suite inside
the loop, not after it. Treat rules whose remediation changes semantics as a
separate class requiring stronger evidence before auto-application. And measure
the gate you are optimising against with an instrument it does not control.

% =====================================================================
\section{Limitations}

\begin{itemize}[leftmargin=1.2em, itemsep=0.35em, topsep=0.4em]
\item \textbf{The breakage result rests on one rule.} \code{S311} supplies 91 of
the 122 security-shown tasks, and without it the paired comparison is $p = 0.09$
on 13 discordant pairs. The \code{S311} finding itself is well powered; a claim
about security rules in general is not, and is not made. \code{B006}'s 53\% is 17
tasks and should be read as a lead.

\item \textbf{Source text was not retained.} Steps record findings, codes, lines
and test outcomes, but not each revision's text, so the mechanism is established
from failure modes and introduced codes rather than from diffs. This is the first
thing to change in a follow-up.

\item \textbf{One sample per task at temperature 0.} The variance budget went into
tasks rather than seeds. These are point estimates for greedy decoding.

\item \textbf{The Pylint arm reports no transfer rate.} Pylint's message ids have
no Ruff or Bandit counterpart, so that arm can say what was addressed but not
whether the fix transferred. It is recorded as unmeasurable rather than as zero
suppression.

\item \textbf{Single-file Python from one benchmark.} Nothing here extends to
Java, to CodeQL, or to repository-scale change without being measured there.
\end{itemize}

% =====================================================================
\section{Reproducibility}

The complete raw record is published as \code{steps.jsonl}: every step of every
arm, with findings by tool and code, test outcome, directives written, and tokens.
All analysis re-derives from it offline, so the tables can be re-sliced or
corrected without another GPU hour. A verification script recomputes all 23 figures
quoted in this paper from that file and exits non-zero if the data stops agreeing;
it was written after an earlier draft's central claim failed its own robustness
check, and it is what caught it.

The run notebook measures the accelerator and sizes itself to it---one A100
($\approx$1.1\,h) or a pair of T4s ($\approx$2.5\,h)---and records the git commit
in its manifest, so a result traces to the code that produced it.

\vspace{0.6em}
{\small\color{muted}
Code and data: \url{https://github.com/SyedMohammedSameer/AgentEval}}

% =====================================================================
\section{Prior work this does not repeat}

Static analysis as a repair feedback loop is established, as is iterative
refinement degrading security, and self-repair across model scales and families.
Patch overfitting to tests is long established in the automated program repair
literature. What appears to remain open, and what this measures, is whether a fix
aimed at one analyser survives contact with another---and what it costs the
program when it does.

\begin{thebibliography}{9}
\small
\bibitem{bcb} Zheng, T. et al.
\emph{BigCodeBench: Benchmarking Code Generation with Diverse Function Calls and
Complex Instructions.} arXiv:2406.15877, 2024.
\bibitem{sa1} \emph{Static analysis feedback in LLM repair loops.} arXiv:2508.14419.
\bibitem{sa2} \emph{Iterative refinement and code security.} arXiv:2412.14841.
\bibitem{sr} \emph{Self-repair across model scales and families.} arXiv:2604.10508.
\end{thebibliography}

\end{document}
